%%%%%%%%%%%%%%%%%%%%%%%%%%%%%%%%%%%%%%%%%%%%%%%%%%%%%%%%%%%%%%%%%%%%%%%%%%%%%%%
\documentclass[a4paper, 11pt]{article}

\usepackage[utf8]{inputenc}
\usepackage[T1]{fontenc}
\usepackage[czech]{babel}
\usepackage[a4paper, margin=1.8cm]{geometry}
\usepackage{amsmath, amsfonts, amssymb, mathtools, bm}
\usepackage{graphicx, hyperref, xcolor, enumitem, booktabs, listings}

\lstset{
    basicstyle=\ttfamily\small,
    backgroundcolor=\color{gray!10},
    frame=single,
    breaklines=true,
    language=Python
}

% Prostředí pro tabulky funkcí
\newenvironment{funkcedict}
{
\renewcommand{\arraystretch}{1.4}
\begin{tabular}{p{0.30\textwidth} p{0.32\textwidth} p{0.33\textwidth}}
\toprule
\textbf{Funkce / Operace} & \textbf{Příklad použití} & \textbf{Vysvětlení} \\
\midrule
}
{
\bottomrule
\end{tabular}
}

\begin{document}

\begin{center}
\Huge \textbf{Vzorcovník -- 4ST204 (Python)}
\end{center}

%% --- SEKCE: TEORIE ----------------------------------------------------------
\section{Statistická teorie}

\subsection*{Míry úrovně a polohy}
\begin{itemize}[noitemsep]
    \item \textbf{Aritmetický průměr:} $\bar{x} = \frac{\sum x_i}{n}$. Klasický průměr pro data bez extrémních hodnot.
    \item \textbf{Geometrický průměr:} $\bar{x}_{G} = \sqrt[n]{\prod x_i}$. Pro multiplikativní data (tempa růstu, úroky).
    \item \textbf{Harmonický průměr:} $\bar{x}_h = \frac{n}{\sum \frac{1}{x_i}}$. Používá se pro poměrové veličiny (např. rychlost v km/h), kde je konstantní čitatel (např. ujetá vzdálenost).
    \item \textbf{Modus ($\hat{x}$):} Nejpravděpodobnější (nejčetnější) hodnota.
    \item \textbf{Kvartily:} Dolní $Q_1$ (25 \%), Medián $Q_2$ (50 \%) a Horní $Q_3$ (75 \%).
    \item \textbf{Sturgesovo pravidlo:} $k = 1 + \log_2 n$ (doporučený počet intervalů).
\end{itemize}

\subsection*{Kvantily (Pozice a výpočet)}
Pozice $p$-tého kvantilu: $z_p = n \cdot p$. Hodnota $\tilde{x}_p$ se určí z uspořádaného souboru $x_{(1)} \le \dots \le x_{(n)}$:
\begin{equation*}
\tilde{x}_p = 
\begin{cases} 
x_{(\lceil z_p \rceil)} & \text{pokud } z_p \notin \mathbb{Z} \\
\frac{x_{(z_p)} + x_{(z_p+1)}}{2} & \text{pokud } z_p \in \mathbb{Z}
\end{cases}
\end{equation*}

\subsection*{$L^p$ Normy (Vzdálenosti a penalizace)}
Obecně: $\|x\|_p = \left( \sum |x_i|^p \right)^{1/p}$. Klíčové pro metriky chyb a regularizaci.
\begin{itemize}
    \item \textbf{$L^1$ (Manhattan):} $\|x\|_1 = \sum |x_i|$. Využití: \textbf{Lasso regrese}.
    \item \textbf{$L^2$ (Euklid):} $\|x\|_2 = \sqrt{\sum x_i^2}$. Využití: \textbf{Ridge regrese} (metoda nejmenších čtverců).
    \item \textbf{$L^\infty$ (Max):} $\|x\|_\infty = \max(|x_i|)$. Reprezentuje nejhorší možnou chybu.
\end{itemize}

%% --- SEKCE: PRÁCE S POLEMI --------------------------------------------------
\section{Práce s poli (Standardní Python)}
\begin{funkcedict}
\textbf{List Comprehension} & \texttt{[x**2 for x in range(5)]} & Dynamické generování pole pomocí cyklu. \\
\texttt{x[start:stop:krok]} & \texttt{pole[1:10:2]} & Výřez (stop je exkluzivní). \texttt{[::-1]} pro otočení. \\
\texttt{x.copy()} vs \texttt{deepcopy} & \texttt{a = b.copy()} & \texttt{copy} je plytká, \texttt{deepcopy} kopíruje i vnořená pole. \\
\texttt{x.append(v)} & \texttt{data.append(10)} & Přidání jednoho prvku na konec pole. \\
\texttt{x.extend(a)} & \texttt{data.extend([1,2,3])} & Přidání více prvků na konec pole. \\
\end{funkcedict}

%% --- SEKCE: ZÁKLADNÍ PYTHON -------------------------------------------------
\section{Základní Python (Built-in \& Math)}
\begin{funkcedict}
\texttt{help(funkce)} & \texttt{help(print)} & Dokumentace pro zadanou funkci. \\
\texttt{range(start, stop)} & \texttt{list(range(0, 5))} & Vygeneruje sekvenci \texttt{[0, 1, 2, 3, 4]}. \\
\texttt{math.log(x, b)} & \texttt{math.log(8, 2)} & Logaritmus čísla $x$ o základu $b$. \\
\texttt{math.sqrt(x)} & \texttt{math.sqrt(16)} & Druhá odmocnina čísla. \\
\texttt{math.ceil / floor} & \texttt{math.ceil(2.1)} & Zaokrouhlení nahoru (3) / dolů (2). \\
\texttt{math.trunc(x)} & \texttt{math.trunc(2.9)} & Odříznutí desetinné části (2). \\
\texttt{math.comb(n, k)} & \texttt{math.comb(10, 3)} & Kombinační číslo (počet kombinací $k$ prvků z $n$). \\
\end{funkcedict}

%% --- SEKCE: NUMPY -----------------------------------------------------------
\section{Knihovna NumPy (Statistika a Matice)}
\begin{funkcedict}
\texttt{np.array([...])} & \texttt{np.array([1, 2, 3])} & Vytvoření základního NumPy pole (vektoru nebo matice). \\
\texttt{np.sum(x)} & \texttt{np.sum(data, axis=0)} & Součet všech prvků v poli (nebo podél zvolené osy). \\
\texttt{np.mean(x)} & \texttt{np.mean(data)} & Průměr všech prvků v poli. \\
\texttt{np.prod(x)} & \texttt{np.prod(data)} & Součin všech prvků v poli. \\
\texttt{np.sqrt(x)} & \texttt{np.sqrt(data)} & Druhá odmocnina z každého prvku v poli. \\
\texttt{np.random.rand(n)} & \texttt{np.random.rand(100)} & Generování pole $n$ náhodných čísel z rovnoměrného rozdělení na $[0, 1)$. \\
\texttt{np.random.uniform(a, b, n)} & \texttt{np.random.uniform(10, 50, 100)} & Generuje $n$ náhodných \textbf{desetinných} čísel v rozsahu $[a, b)$. \\
\texttt{np.quantile(x, q)} & \texttt{np.quantile(x, 0.75)} & Výpočet kvantilu ($q \in [0, 1]$). \\
\texttt{np.log10(x)} & \texttt{np.log10(data)} & Dekadický logaritmus (základ 10). Často se používá pro normalizaci dat. \\
\texttt{np.std(x, ddof=1)} & \texttt{np.std(data, ddof=1)} & Výběrová směrodatná odchylka (\texttt{ddof=1} je nutné pro nestranný odhad). \\
\texttt{np.linalg.norm(x, p)} & \texttt{np.linalg.norm(x, ord=2)} & Výpočet $L^1, L^2$ nebo $L^\infty$ (\texttt{np.inf}) normy. \\
\texttt{np.insert / np.delete} & \texttt{np.insert(A, 1, [..], axis=0)} & Vloží/Smaže řádek (\texttt{axis=0}) nebo sloupec 
(\texttt{axis=1}). \\
\texttt{stats.gmean(x)} & \texttt{stats.gmean(x)} & Geometrický průměr (knihovna Scipy). \\
\texttt{Geometrický průměr} & \texttt{10**np.mean(np.log10(x))} & Výpočet geometrického průměru pomocí logaritmů. \\
\texttt{stats.hmean(x)} & \texttt{stats.hmean(x)} & Harmonický průměr (knihovna Scipy). \\
\end{funkcedict}

%% --- SEKCE: PANDAS ----------------------------------------------------------
\section{Knihovna Pandas (Analýza dat)}
\begin{funkcedict}
\texttt{df.describe()} & \texttt{df.describe()} & Souhrnné statistiky (mean, std, kvartily, atd.). \\
\texttt{df.loc} vs \texttt{df.iloc} & \texttt{df.loc[0, 'Sloupec']} & Výběr dat podle \textbf{názvů} (\texttt{loc}) vs. \textbf{pozice} (\texttt{iloc}). \\
\texttt{df.value\_counts()} & \texttt{df['A'].value\_counts()} & Tabulka četností pro kategorická data. \\
\texttt{df.isna().sum()} & \texttt{df.isna().sum()} & Počet chybějících (null) hodnot v tabulce. \\
\end{funkcedict}

\end{document}
